\section{Discussion}
\label{sec:discussion}

\subsection{BoundedField}
BoundedField is a cheap performance uplift over any embedder, as it does not involve LLM training or invocation, and its scoring runs entirely on the CPU. It only needs a citation graph, which GraphIV already ships. At the cutoffs of \cref{fig:embedders}, scaling the embedder from Nomic's 475M parameters to Harrier's 27B raised cosine recall by 0.073 on SurveyMatch, 0.028 on FrontierMatch, and 0.151 on ConsensusMatch, while BoundedField over Nomic alone raised it by 0.187, 0.072, and 0.452. The two gains add up, as BoundedField still improved Harrier 27B's recall by 0.157, 0.057, and 0.354. This can be attributed to how embedders provide semantic context while BoundedField enables citation-driven scores to be considered. Neither is exclusive to the other.

The recipes are also largely embedding-agnostic. The reported recipes were trained on Nomic and applied unchanged to the other embedders, yet they still improved over cosine in 50 of the 64 pairs (\cref{tab:transfer}), and in every pair outside RealScholarQuery and ScholarCatalyst. All 14 losses were on RealScholarQuery and ScholarCatalyst, whose queries ask for specific papers that a stronger embedder already ranks high, and refitting the \term{focused} recipe on each embedder recovered six of them.

\term{anc}, \term{cocite}, and the citation hop are the most impactful for BoundedField (\cref{tab:ablation}), as what a review cites is largely what the papers most similar to the question cite, which can be recovered via citation graph traversal, not semantic similarity. This also explains why FrontierMatch was hardest, as its answers are recent papers the core had little time to cite, and \term{anc} is zero for an uncited paper.

\subsection{Dataset and Benchmarks}
Apart from the topic tree, each stage of GraphIV processes only incremental updates, meaning biweekly snapshots add papers and citations without a full rebuild, allowing us to keep BoundedField's graph current. The hash split of the benchmarks only adds cases, keeping earlier cases in place. This update frequency and the intent labels set GraphIV apart from earlier open graphs. The labels let agents built on GraphIV determine which papers extended or critiqued a given work, which is impossible with unlabeled citation edges. As every answer is a paper a later review cited, the benchmarks reward work already being noticed, and a paper too new to be cited cannot be an answer.

\subsection{Automated Research System}
Every search the agent ran was a BoundedField search, and we found it picked modes freely (\cref{tab:modes}), never picking \term{consensus}, as no question asked for it. The agent was able to compose a chain of calls (\cref{tab:runs}) that led to the final report, and \cref{app:trace} follows one run call by call. The mechanical check kept the reports grounded: 61 of the 63 first drafts failed it on statements with no source, numbers absent from their paragraph, or names the cited papers never mention (\cref{tab:faults}). As the check matches strings, it guarantees by construction, whatever the model, that every quote exists in its paper and every cited paper was read, but not that a statement fully agrees with the sentence it cites.

\subsection{Cost}
This project was created on a single RTX 3090. Labeling the citations cost 12 GPU hours, and embedding the corpus cost 5. Biweekly updates are network-bound and add hardly any GPU hours. The release, omitting the citing sentences, is 6.7 GB of Parquet. BoundedField took about one second per query against 0.67 seconds for PPR on the same pool and 0.41 for cosine search, and a \term{standard} report took 12.7 minutes on average, with no paid API call.

\subsection{Limitations and Future Work}
GraphIV can be scaled in volume, as it currently covers only arXiv and parses only papers with an HTML render. The citation intents, review scores, and topic names were produced by LLMs or classifiers distilled from one, with a rather small portion (985 cases) labeled by a human. The errors accumulated in the process are likely to propagate into the lineages. Gemma 4 26B also wrote the benchmark questions, about 6\% of the papers the review classifier accepts at our threshold are not reviews, and the numbers reported by others in \cref{tab:existing,tab:scholarcatalyst} were not rerun. The agent was tested on 21 questions, once per effort, with one model, and no human expert rated the reports. 

We plan to grow GraphIV along three axes. In \emph{volume}, it would extend past arXiv in tiers, keeping arXiv, whose full text is open, at the most detail while adding other sources at the depth their access enables, potentially resolving most of the remaining citations. In \emph{labels}, human annotators would label more citations, drawing on diverse opinions, and give feedback on the final automated literature reviews. Finally, in \emph{resolution}, we plan to embed every paragraph of open papers rather than the abstract alone, providing a claim- and paragraph-level citation graph to begin the creation of a high-resolution graph of science.
